\documentclass[10pt]{article}
%%%%%%%%%%%%%
%windows 正常 ok
\usepackage{ctex} % 加载中文支持宏包
%macos ok
%\usepackage{fontspec}
%\usepackage{xeCJK}  % 如果你需要支持中文
%%%%%%%%%%%%%

\usepackage[x11names]{xcolor} % 启用更多颜色选项
\usepackage{amsmath}      % 数学公式
\usepackage{graphicx}     % 图片插入
\usepackage{listings}     % 用于插入代码
\usepackage{hyperref}     % 超链接支持
\usepackage{geometry}     % 设置页边距
\usepackage{fancyhdr}     % 页眉页脚设置
\usepackage{multicol}   % 多栏排版
\usepackage{tocbibind}    % 目录、参考文献包含在目录中
\geometry{a4paper, left=0.6in, right=0.5in, top=0.5in, bottom=0.5in}

% 设置列表中Java代码的样式
\lstdefinestyle{java}{
    language=Java,
%    backgroundcolor=\color{lightgray}, % 背景颜色
    basicstyle=\ttfamily\footnotesize, % 字体
    breaklines=true,                   % 自动换行
    captionpos=b,                      % 标题位置
%    numbers=left,                      % 行号位置
%    numberstyle=\tiny\color{gray},     % 行号样式
    keywordstyle=\color{blue},         % 关键词颜色
    stringstyle=\color{red},           % 字符串颜色
    commentstyle=\color{green},        % 注释颜色
%    frame=single                       % 边框
}

% 设置页眉和页脚
\pagestyle{fancy}
\fancyhf{}
\fancyhead[L]{技术学习} % 左边页眉
\fancyhead[C]{sully} % 中间页眉
\fancyhead[R]{202501} % 右边页眉
\fancyfoot[C]{\thepage} % 页脚中间显示页码

\title{技术学习}
\author{sully}
\date{\today}

\begin{document}

\maketitle

% 生成目录
\tableofcontents
\newpage

\section{引言}
这里是文章的引言部分。

\section{Java代码示例}
以下是一个嵌入的Java代码示例：

\begin{lstlisting}[style=java, caption={HelloWorld.java}]
public class HelloWorld {
    public static void main(String[] args) {
        System.out.println("Hello, world!");
    }
}
\end{lstlisting}

\section{图片插入示例}
下面是一个插入图片的示例：

\begin{figure}[h]
    \centering
    \includegraphics[width=0.5\textwidth]{9.png} % 确保替换为你的图片路径
    \caption{示例图片}
    \label{fig:example}
\end{figure}


\newpage
\section{引言}
这里是文章的引言部分。

\newpage % two sum 不同语言解法，含冒泡和插入排序
%\section{Java代码示例}
%\vspace{-5mm}
% 开始多栏排版
\begin{multicols}{2}
%\lstset{aboveskip=-2mm}  % 控制 lstlisting 上方的空白
\begin{lstlisting}[style=java, caption={}]

class Solution {
    public int[] twoSum(int[] nums, int target) {
        int[] aaa = { 7, 5, 3, 2, 8, 0, 1 };

        System.out.println("冒泡排序");
        for (int i = 0; i < aaa.length; i++) {
            for (int j = 0; j < aaa.length - i - 1; j++) {
                int tmp = aaa[j + 1];
                if (aaa[j] > aaa[j + 1]) {
                    aaa[j + 1] = aaa[j];
                    aaa[j] = tmp;
                }
            }
        }

        for (int i = 0; i < aaa.length; i++) {
            System.out.print(aaa[i] + ",");
        }

        System.out.println("");
        System.out.println("插入排序");
        int[] bbb = { 7, 5, 3, 2, 8, 0, 1 };
        for (int i = 1; i < bbb.length; i++) {
            int tmpV = bbb[i];
            int j = i - 1;

            /* 方法一，将比 key 大的元素向右移动 */
            /*
            while (j >= 0 && bbb[j] > tmpV) {
                bbb[j + 1] = bbb[j];
                j = j - 1;
            }
            bbb[j + 1] = tmpV;
            */
            //方法二
            for(;j>=0;j--) {
                if(bbb[j]>tmpV){
                    bbb[j+1] = bbb[j];
                } else {
                    break;
                }
            }
            bbb[j+1] = tmpV;

        }
        for (int i = 0; i < bbb.length; i++) {
            System.out.print(bbb[i] + ",");
        }

        System.out.println("two sum");
        Map <Integer,Integer> myMap = new HashMap<>();
        for(int i=0;i<nums.length;i++){
            if(myMap.containsKey(target-nums[i])) {
                return new int []{myMap.get(target-nums[i]),i};
            }
            myMap.put(nums[i],i);
        }
        return new int [0];
    }
  }

def twoSum(self, nums: List[int], target: int) -> List[int]:
        n = len(nums)
        for i in range(n - 1):
            for j in range(i + 1, n):
                if nums[i] + nums[j] == target:
                    return [i, j]
        return []  # No solution found


class Solution:
    def twoSum(self, nums: List[int], target: int) -> List[int]:
        numMap = {}
        n = len(nums)

        for i in range(n):
            complement = target - nums[i]
            if complement in numMap:
                return [numMap[complement], i]
            numMap[nums[i]] = i

        return []  # No solution found
go

func twoSum(nums []int, target int) []int {
    // Hash table to store number->index mapping
    ht := make(map[int]int)

    // Iterate through the array
    for i, num := range nums {
        // Check if complement exists in hash table
        if j, exists := ht[target-num]; exists {
            // If found, return indices of both numbers
            return []int{j, i}
        }

        // Store current number and its index in hash table
        ht[num] = i
    }

    // Return empty slice if no solution found
    return []int{}
}
javascript

var twoSum = function(nums, target) {
    let map = new Map();

    for (let i=0; i<nums.length; i++){
        if(map.has(target-nums[i])){
            return [map.get(target-nums[i]), i]
        }else{
            map.set(nums[i], i)
        }
    }
};

\end{lstlisting}

%\columnbreak % 强制换栏，继续在第二栏排版
%\begin{lstlisting}[style=java, caption={}]
%public class HelloWorld {
%  public static void main(String[] args) {
%        // 你好，世界
%        System.out.println("Hello, world!");
%    }
%}
%\end{lstlisting}
\end{multicols}

\newpage
% 开始多栏排版
\begin{multicols}{2}
\begin{lstlisting}[style=java, caption={}]
//c++
class Solution {
public:
    vector<int> twoSum(vector<int>& nums, int target) {
        unordered_map<int, int> numMap;
        int n = nums.size();

        for (int i = 0; i < n; i++) {
            int complement = target - nums[i];
            if (numMap.count(complement)) {
                return {numMap[complement], i};
            }
            numMap[nums[i]] = i;
        }

        return {}; // No solution found
    }
};

N.283 移动零问题
public void moveZeroes(int[] nums) {
        int j =0;
        for (int i = 0; i < nums.length; i++) {
            if (nums[i]!=0) {
                nums[j] = nums[i];
                if (i!=j){
                    nums[i] =0;
                }
                j++;
            }
        }}
\end{lstlisting}
\end{multicols}

\newpage
\section{插入C++代码}
%\lstinputlisting[language=java]{Hello.java}
%\input{Hello.txt}
%\usepackage{verbatim}
%\verbatiminput{Hello.txt}

本文展示了如何嵌入Java代码和图片，同时提供了目录、章节标题等结构。

\end{document}
